\documentclass[11pt,a4paper]{article}
\usepackage[margin=26mm]{geometry}
\usepackage{amsmath,amssymb,amsthm,hyperref}
\newtheorem{lemma}{Lemma}\newtheorem{proposition}[lemma]{Proposition}
\newcommand{\F}{\mathbb F}\newcommand{\nuu}{\nu_2}
\newcommand{\PCO}{\operatorname{PCO}}
\newcommand{\Om}{\Omega}
\setlength{\emergencystretch}{2em}
\title{Bounds for orthogonal fixed groups in odd characteristic}
\author{Human Compression project}\date{}
\begin{document}\maketitle
This component supplies fixed-group bounds and two subspace counts for the orthogonal generation argument. Throughout $q$ is odd and
$\nuu(H)=|\{h\in H:h^2=1\}|$, including the identity.
We write $\PCO(V)$ for the full conformal orthogonal group modulo
the scalar group over its ground field. This avoids confusion with
the image of the isometry group, which can be a proper subgroup.

We retain the standard finite classical group orders, the faithful action
of projective conformal groups on their derived simple groups, and
Lawther--Liebeck--Seitz, Proposition~1.3~\cite{LLS}. The latter gives
$\mathrm i_2(\operatorname{Aut}S)<2(Q^{N+r}+Q^{N+r-1})$ over $\F_Q$.
The actual torus centralizers and projective eigenvalue degrees are
acquired from the separately reviewed orthogonal torus component~\cite{tori}.

\begin{lemma}\label{small}
For $d\ge2$ and a nondegenerate $d$-space over $\F_q$, put
$e_d=\lfloor d^2/4\rfloor$. Then
\[
 \nuu(\PCO_d(q))\le3q^{e_d},\qquad
 |\PCO_d(q)|<3q^{d(d-1)/2}.
\]
\end{lemma}
\begin{proof}
For $d\ge5$, the derived projective group is simple, including the
identifications in dimensions five and six. Its exponent $N+r$ is $e_d$.
The faithful projective action and the retained bound give
$\nuu(\PCO_d(q))<2q^{e_d}+2q^{e_d-1}+1\le3q^{e_d}$.

In dimension two the group is dihedral of order $2(q-\epsilon)$;
its number of elements of square one is $q-\epsilon+2\le q+3$.
In dimension three it is $\operatorname{PGL}_2(q)$, with $q^2$
nonidentity involutions. In dimension four the two models, already
verified in the linear/unitary component~\cite{linear}, are
\[
 \operatorname{PGL}_2(q)\wr C_2,\qquad
 \operatorname{PGL}_2(q^2)\rtimes C_2.
\]
The respective values of $\nuu$ are
$q^4+q^3+2q^2-q+1$ and $q^4+q^3+q+1$, both at most $3q^4$.
This proves the first assertion in every dimension needed here.

For odd $d$, the projective conformal group has the order of
$\operatorname{SO}_d(q)$, less than $q^{d(d-1)/2}$.
For even $d=2s$, it has the order of the full isometry group,
\[
 2q^{s(s-1)}(q^s-\epsilon)\prod_{j=1}^{s-1}(q^{2j}-1)
 <2(1+q^{-s})q^{d(d-1)/2}\le\tfrac83 q^{d(d-1)/2}.
\]
The ordinary order formulas also include dimension two.
\end{proof}

\begin{proposition}\label{balanced}
Let $V$ have dimension $2m$, $m\ge3$. Suppose a conformal map $a$
has square $cI$, with multiplier $c$, and its two eigenspaces over the
algebraic closure both have dimension $m$. Then
\[
 \nuu(C_{\PCO(V)}([a]))\le24q^{2\lfloor m^2/4\rfloor}.
\]
If $c$ is a square and the two eigenspaces instead have unequal
dimensions $r,s\ge2$, then
\[
 \nuu(C_{\PCO(V)}([a]))\le18q^{e_r+e_s}.
\]
\end{proposition}
\begin{proof}
First suppose $c$ is a square. Rescale $a$ to have square one.
Its two eigenspaces are nondegenerate and orthogonal. A map centralizing
$[a]$ sends $a$ to $a$ or $-a$, since it preserves the multiplier.
It therefore preserves the two eigenspaces, or interchanges them.
Projecting its restrictions to their projective conformal groups has
kernel of order at most two: a scalar on each block must have the same
square, because the two multipliers agree, and a common scalar is
already factored out.

When an interchange is possible the two forms are similar. Writing
$H$ for either projective conformal group, the image is a subgroup of
$H\wr C_2$. There are $\nuu(H)^2$ elements of square one in the base,
and exactly $|H|$ in the interchange coset, since its two coordinates
must be mutual inverses. A group extension with kernel of order at
most two has at most twice as many elements of square one as its image.
Lemma~\ref{small} thus gives
\[
 \nuu(C)\le2\bigl(9q^{2e_m}+3q^{m(m-1)/2}\bigr)
 \le24q^{2e_m}.
\]
If the forms are not similar, omit the interchange coset and use the
product bound. Unequal dimensions also forbid interchange, giving
$2(3q^{e_r})(3q^{e_s})$ as stated.

Now suppose $c$ is a nonsquare. The algebra $K=\F_q[a]$ is $\F_{q^2}$.
The map $a$ is self-adjoint for the underlying bilinear form, because
$a^*=ca^{-1}=a$. The nondegenerate trace pairing then writes that form
uniquely as the trace of a nondegenerate symmetric $K$-bilinear form on
the $m$-dimensional $K$-space $V$. A projective centralizer acts either
linearly or by the nontrivial field automorphism of $K$; the induced
projective group embeds in the automorphism group of the corresponding
derived orthogonal group over $K$. Its kernel has order at most two.
Indeed a scalar $z\in K^\times$ induces an allowed similitude over
$\F_q$ only if $z^2\in\F_q^\times$, and there are exactly two such
cosets modulo $\F_q^\times$.

For $m\ge5$, simplicity and the retained involution bound give
$\nuu(C)\le6q^{2e_m}$. The same argument applies when $m=3$, using
$\operatorname{P}\Omega_3(K)\simeq\operatorname{PSL}_2(K)$.
For $m=4$ of minus type it uses
$\operatorname{P}\Omega_4^-(K)\simeq\operatorname{PSL}_2(q^4)$,
and gives $6q^8$.

The plus case in dimension four needs its two simple factors. Put
$Q=q^2$ and $A=\operatorname{PSL}_2(Q)$; the derived projective group
is $A\times A$. Its full automorphism group is
$\operatorname{Aut}(A)\wr C_2$. We have
$\nuu(\operatorname{Aut}A)\le3Q^2$ and
$|\operatorname{Aut}A|=f_Q Q(Q^2-1)<Q^4$, where
$Q=p^{f_Q}$ and $f_Q\le Q$. Thus
\[
 \nuu(C)\le2\bigl(9Q^4+Q^4\bigr)=20q^8.
\]
These bounds all imply the stated constant $24$.
\end{proof}

\begin{proposition}\label{application}
For a natural involutory reversing automorphism of any of the chosen
orthogonal elements over odd $\F_q$, its fixed subgroup in the simple
group has at most
\[
 \begin{cases}
 24q^{2\lfloor m^2/4\rfloor},&\dim V=2m,\\
 18q^{m(m+1)/2},&\dim V=2m+1
 \end{cases}
\]
elements of square one. This assertion does not count triality
reversers which lack a representative on the natural module.
\end{proposition}
\begin{proof}
Let $X$ be the natural representative of the chosen torus element.
Projective inversion initially gives $aXa^{-1}=\lambda X^{-1}$,
with $\lambda^2=1$. On the large block, the acquired degree lemma
forces the field exponent of $a$ to be $k$, where the block has
dimension $2k$. Since the eigenvalue $\alpha$ satisfies
$\alpha^{q^k}=\alpha^{-1}$ exactly, the scalar is $\lambda=1$.
Consequently the restriction is $z\operatorname{Frob}^{k}$.
Its square and conformal multiplier are both
$\Norm_{\F_{q^{2k}}/\F_{q^k}}(z)$, which lies in $\F_q^\times$.
The same calculation on the negative two-plane gives the same scalar;
on a fixed line its multiplier is the square of its line scalar.
Thus $a^2=cI$, with multiplier $c$.

If $c$ is a square, on every negative $2k$-block a scalar multiple of
$a$ is a semilinear involution for the quadratic extension over
$\F_{q^k}$. Its fixed and negative spaces both have dimension $k$.
This follows also by solving $zv^{q^k}=\pm v$: norm one is precisely
the solvability condition, and each solution space is one-dimensional
over $\F_{q^k}$. The dimensions are therefore $m,m$ in even dimension
and $m,m+1$ in odd dimension. In the latter case the line forces $c$
to be a square. Apply Proposition~\ref{balanced}; taking a subgroup
cannot increase its number of elements of square one.
\end{proof}


\begin{proposition}\label{lower}
For odd $q$ the following strict lower bounds hold:
\[
 \mathrm i_2(P\Om_8^+(q))>q^{16}/4,
 \qquad \mathrm i_2(P\Om_8^-(q))>q^{16}/2.
\]
\end{proposition}
\begin{proof}
In either space let $W$ run over the nondegenerate four-dimensional
subspaces of plus type. Put $t_W=-I$ on $W$ and $t_W=I$ on its
orthogonal complement. Its determinant is one and its spinor norm is
the discriminant of the four-dimensional plus form, a square.
Thus $t_W\in\Om(V)$. We use here the standard spinor norm of $-I$
on a nondegenerate subspace, as in~\cite[Definition~3.1]{MRD}.
Witt's extension theorem makes the full isometry group transitive on
these subspaces, with stabilizer $O(W)\times O(W^\perp)$.

In plus dimension eight the complement also has plus type. The projective
images of $t_W$ and $t_{W'}$ agree exactly when $W'=W$ or $W'=W^\perp$:
the scalar identifying two isometries must be $1$ or $-1$.
No nondegenerate $W$ equals its complement. Hence the resulting distinct
involutions number
\[
 \frac{|O_8^+(q)|}{2|O_4^+(q)|^2}
 =\frac{q^8(q^2+1)^2(q^4+q^2+1)}4>\frac{q^{16}}4.
\]
In minus dimension eight the complement has minus type, while $-I$
has nonsquare spinor norm and does not lie in $\Om(V)$.
There is no projective identification between the displayed involutions.
Their number is
\[
 \frac{|O_8^-(q)|}{|O_4^+(q)|\,|O_4^-(q)|}
 =\frac{q^8(q^4+1)(q^4+q^2+1)}2>\frac{q^{16}}2.
\]
The count is of actual subspaces, so it does not require transitivity
of the smaller omega group on this set.
\end{proof}

\paragraph{Scope.}
The constants deliberately include the identity, possible projective
order-four lifts, a nonsplit quadratic-field form, and interchange
cosets. No central quotient is estimated by counting only involutions
in its covering group. The separate covering and total-involution
estimates remain necessary for a generation conclusion.

\begin{thebibliography}{9}\small
\bibitem{LLS} R.~Lawther, M.~W.~Liebeck and G.~M.~Seitz,
\emph{Fixed point ratios in actions of finite exceptional groups of Lie type},
Pacific Journal of Mathematics 205 (2002), Proposition~1.3.
The full automorphism bound, including the identity adjustment here,
is retained, not reproved.
\bibitem{MRD} S.~H.~Murray and C.~M.~Roney-Dougal,
\emph{The spinor norm and homomorphism algorithms for classical groups},
preprint dated 2 May 2008, Definition~3.1 and Lemma~3.2.
\bibitem{tori} Human Compression project,
\emph{Orthogonal tori and their projective orders}, accepted component
with independent review.
\bibitem{linear} Human Compression project,
\emph{Linear and unitary groups}, accepted component; the precise
dimension-four projective conformal models and their counts are acquired.
\end{thebibliography}
\end{document}
